\section{Numerical Solution by Operator Splitting}\label{sec:split}%
\label{sec:why-split}

The system of \cref{sec:summary-equations} on the mesh of \cref{sec:geometry}
carries $2\,446\,659$ immobile and $362\,468$ mobile degrees of freedom, and as
one coupled stiff system it would be hopeless. Three properties compound. It is
stiff: recombination acts on microseconds and the loops evolve over
$10^{7}\,$s, so an implicit integrator is forced and must factorize a Jacobian. It is
large: that factorization is $O(N^{3})$ at $N=2\,809\,127$, and sparsity softens
without rescuing it, a three-dimensional finite-element Jacobian still going as
$\sim O(N^{2})$ with fill as $O(N^{4/3})$. And it is long: a march to 10\,dpa at
$10^{-7}\,\mathrm{dpa\,s^{-1}}$ is $10^{8}\,$s, over which an adaptive
integrator takes $\sim10^{2}$ internal steps per unit of slow dynamics.

What makes the problem hard is also what resolves it. The mobile species reach
quasi-steady state on $\tau_M\sim(k^{2}D)^{-1}$, of order microseconds at
$573\,$K, while the immobile families evolve on the dose scale --- thirteen
decades away. On the slow manifold they are therefore slaved: writing the two
blocks as $\dot{\spvec c}_M=\varepsilon^{-1}\mathcal R_M(\spvec c_M;\spvec Q)$
and $\dot{\spvec Q}=\bm g(\spvec Q;\spvec c_M)$ with
$\varepsilon=\tau_M/\tau_Q\ll1$, the singular limit $\varepsilon\to0$ replaces
the first by the algebraic constraint
$\mathcal R_M(\spvec c^{\star}_M;\spvec Q)=\bm 0$. The spatial code implements
that limit exactly: its mobile solver carries no time derivative, so
\cref{eq:master-c} is the steady problem itself, not a transient taken to
large times.

The two blocks are advanced by a Lie--Trotter split. Over one dose interval,
with $\Delta t=(\gamma_{n+1}-\gamma_n)/G$, the scheme alternates four steps:

\begin{enumerate}[label=\textbf{S\arabic*.},itemsep=2pt]
\item Solve the elastic problem with the current inelastic deformation for
      $\bm\sigma_n(\xv)$ --- constant here, retained for generality.
\item Assemble the local coefficients from $\spvec Q_n(\xv)$ and $\bm\sigma_n$:
      the diffusion tensors, the efficiencies of \cref{eq:Zsk}, the sink
      operator of \cref{eq:KM}, $\spmat R_2$, the variant weights and the
      boundary data.
\item With $\spvec Q_n$ fixed, solve
      $\mathcal R_M(\spvec c^{\star}_M;\spvec Q_n)=\bm 0$ for the quasi-steady
      mobile field --- finite element, Newton, spatially coupled.
\item At every node, integrate \cref{eq:block-n,eq:block-c,eq:block-q} over
      $\Delta t$ with the four mobile values frozen at $\spvec c^{\star}_M$ ---
      pointwise and decoupled.
\end{enumerate}

Two properties of that loop are easily lost, both found by their failure modes.
The fast step must be inside it: replayed instead from a recorded calculation,
$\spvec c_M$ never responds to the immobile state the march is building, so a
state away from quasi-steady state can never relax --- and any spatially uniform
seed is such a state, carrying no boundary layer where the true
$\spvec c^{\star}_M(\xv)$ is pinned to equilibrium on every Dirichlet face. And solver settings must be shared across compared routes: the
spatial code's default makes the mobile solve a projected Newton iteration
capping out at $53.6\%$ error in $c^{i}$, so it is disabled throughout ---
left on one side only, a comparison measures that defect, not the physics.

With $\spvec c^{\star}_M(\xv)$ held fixed the immobile equations at different
nodes are independent --- not an approximation made for speed, but because the
only thing coupling neighboring points is diffusion of the mobile species. The step reuses the same
stiff integrator with a flag that zeroes the four mobile derivatives
\emph{after} every reaction term has been evaluated, so the immobile equations,
the accumulators and $\rho_N$ all read those values as frozen parameters. Each node is one case of an OpenMP batch with
its own workspace, and identical states are deduplicated before dispatch --- a
measured $1.24$ at the end of the march, higher early on when much of the domain
shares one state. The integrator was chosen by measurement:
CVODE with backward differentiation formulas and a dense solver at
$\mathrm{rtol}=10^{-6}$, $\mathrm{atol}=10^{-20}$, whose error against a tight
reference is $3.0\times10^{-4}$, where a Krylov solver gives $1.7\times10^{-2}$
and the Runge--Kutta tables designed as IMEX pairs fail outright in pure
implicit mode. The frozen system is $36$-dimensional, so a dense factorization
is $\sim1.6\times10^{4}$ flops and sparsity offers nothing.

That decoupling makes the cost bearable. The slow step integrates
$3\,262\,212$ scalar ordinary differential equations per substep --- but not as
one system of three million. It is $72\,928$ independent systems of $36$, and
the distinction is worth ten orders of magnitude:

\begin{center}
\begin{tabular}{@{}lr@{}}
\toprule
one implicit factorization & flops \\
\midrule
split, $159\,271\times36^{3}$ & $7.4\times10^{9}$ \\
coupled, $(189\,533\times31)^{3}$ & $2.0\times10^{20}$ \\
\midrule
ratio & $2.7\times10^{10}$ \\
\bottomrule
\end{tabular}
\end{center}

Stiffness localizes the same way: each point's system genuinely is stiff, of
order $100$ internal steps per substep, which is why BDF, but its cost scales
with the dimension of the coupled block, not with how many blocks there are. The coupling did not vanish, then; it moved to where it is affordable. The
tightly coupled problem is now the fast solve, $758\,132$ unknowns and
$3\,\mathrm{m}\,00\,\mathrm{s}$ on 40 cores against
$52.7\,\mathrm{s}$ for an entire slow sweep over all $189\,533$
points, and what the split buys is frequency: with no time derivative in
\cref{eq:master-c} it is solved 31 times over the march rather than every
timestep. Taking the five Newton iterations per solve measured on the smaller
case, hence $\sim36\,$s per linear solve, the counterfactual is
$100\times160\div5\approx3\,200$ factorizations, or $32\,$h for the mobile
block \emph{alone}, against the $2\,\mathrm{h}\,21\,\mathrm{m}$ the slow side
actually took --- and it is linear in node count, as the structure predicts.

The bill arrives as splitting error. Freezing $\spvec c^{\star}_M$ over a
substep ignores the drift of $\spvec Q$, and hence of the sinks, within it --- a
first-order Lie--Trotter error, $O(\Delta\gamma)$ per step and one-sided: with
strengthening sinks the loops are slightly over-grown. The accuracy dial is
therefore the fast-solve cadence, not the dose-snapshot spacing --- the march
takes $n_{\mathrm{sub}}$ substeps per interval and re-solves the mobile field
every $n_{\mathrm{fem}}$ of them. Two consequences follow, again
from failure modes: the cadence counter must run over the whole march, since
counted within an interval it goes silently inoperative whenever that interval
holds fewer substeps than the cadence, and the substep count is fixed per dose
interval rather than by a wall-clock cap on $\Delta t$, a 5\,dpa interval here
being $5\times10^{7}\,$s against which a $5\times10^{4}\,$s cap would have
spawned a thousand subprocesses. Three refinements go unused: a step cap on
$\max_q\Delta\langle k^{2}\rangle/\langle k^{2}\rangle$; Strang splitting,
$O(\Delta\gamma^{2})$ for one extra fast solve; and an in-step Picard iteration
between S3 and S4, which removes the error entirely.

\Cref{tab:cost} reports the measured cost on 40 cores, and two readings matter.
The slow side now dominates, $57\%$ against $38\%$, because this march re-solves
the mobile field only twice per dose interval
($n_{\mathrm{sub}}=10$, $n_{\mathrm{fem}}=5$) where the smaller case re-solved
it every third substep; raising the cadence buys accuracy at nearly linear cost
in the fast side alone, and there is room to do so here. And one
point-integration costs $331\,\mu\mathrm s$ for $36$ stiff equations over a dose
interval, which is what makes twenty-five million of them over the march four
hours rather than a month.

\begin{table}[htbp]
\caption{Measured cost of the 1000\,nm hexagonal single-crystal march of
\cref{sec:results}, on 40 cores. The state carried is 40 wide --- 4 mobile,
27 immobile, 6 conservation accumulators, the network density, and the two
grain-boundary accumulators of \cref{sec:gb-loops} --- of which 36 are
integrated, the four mobile values being held fixed over each substep.}
\label{tab:cost}
\centering\small
\begin{tabular}{@{}lr@{}}
\toprule
Quantity & Value \\
\midrule
CD nodes & 189\,533 \\
mobile field, 4 species & 758\,132 dof \\
immobile field, 9 families $\times$ 3 moments & 5\,117\,391 dof \\
elastic displacement & 568\,599 dof \\
\midrule
state carried per point & 40 \\
integrated per point (= Newton block) & 36 \\
substeps & 160 \\
points integrated per substep, after deduplication & 159\,271 of 189\,533 \\
deduplication factor & $\times1.19$ \\
point-integrations over the march & 25\,483\,360 \\
scalar ODEs integrated over the march & $9.2\times10^{8}$ \\
\midrule
fast solves ($31\times$ finite-element Newton) & 1\,h\,32\,m\,52\,s \quad(37.8\%) \\
slow substeps ($160\times$ batch CVODE) & 2\,h\,20\,m\,38\,s \quad(57.3\%) \\
snapshot and checkpoint writes & 11\,m\,54\,s \quad(4.8\%) \\
\textbf{march total} & \textbf{4\,h\,05\,m\,24\,s} \\
\midrule
one fast solve & 3\,m\,00\,s \\
one slow substep & 52.7\,s \\
one point-integration & 331\,$\mu$s \\
\bottomrule
\end{tabular}
\end{table}

\paragraph{The costs above are measured after an optimization pass, and none of
it is an approximation.} A march of these sixteen intervals on this domain and
mesh ran $3.14\times$ slower before it --- \SI{15.36}{\hour} against
\SI{4.89}{\hour} --- and the acceptance test throughout was bit-identity of the
saved march state, not agreement to a tolerance. Five changes account for it.
The slow batch is \emph{sharded} across processes rather than threaded inside
one, which matters because the threading did not scale at all: at fixed work,
eighty threads in a single process ran slower than one thread, while the same
eighty arranged as twenty processes of four ran $13.2\times$ faster than serial.
The per-case solver input is no longer materialized token by token, which had
consumed $47\%$ of the slow step in string handling for input that is the same
$\sim$87-token base with the state slots overwritten. The finite-element
sparsity pattern is recorded once and reused across Newton iterations, since
only the values move between them. The fast step's thread count is set
explicitly, Windows otherwise confining a process to one of its two
forty-processor groups. And the finite-element program is held \emph{resident}
across the march, so that the Cholesky factorization of the diffusion operator
--- \SI{176}{\second} of every \SI{399}{\second} invocation at this size, and
growing as $N^{1.86}$ against the solve's own $N^{1.33}$ --- is performed once
rather than at each fast solve. That last change is the one that grows with the
problem: it takes a single fast solve from \SI{510}{\second} to
\SI{200}{\second} here, while on a \SI{200}{\nano\meter} case the same rebuild
is under half a minute and the gain correspondingly smaller. The method is
documented separately~\citep{dislocluster_opt}; nothing in it changes an
equation, a tolerance or a discretization.
